\section{Introduction}
\label{sec:intro}

Transparency logs, such as Certificate Transparency (CT)~\cite{rfc6962},
have become a cornerstone of verifiable accountability for critical digital
infrastructure. Recent work has extended this paradigm to Artificial
Intelligence: BTV-Transparency~\cite{btv-t-repo} demonstrates that
high-stakes AI decisions can be structurally bound to public, append-only
logs via linear resource types, guaranteeing that no decision can be
delivered without a verifiable persistence record. This architecture
satisfies the auditability requirements of the EU AI Act Art.~86 and
creates a tamper-evident audit trail for downstream regulators.

However, this immutable architecture conflicts directly with data
protection rights. GDPR Art.~17 and LGPD Art.~18 grant individuals the
right to erasure of their personal data. While this right is fundamental
for privacy, it introduces a novel attack surface against algorithmic
accountability: the \emph{Selective Deletion Attack}.

Unlike classical adversaries who seek to steal data or disrupt service,
a system operator facing regulatory scrutiny for discriminatory practices
has a strong financial incentive to sanitise their historical records.
Enforcement actions --- such as the SEC record-keeping fines of
2024~\cite{sec-fines-2024}, where the absence of records precluded
verification of misconduct --- illustrate the value operators place on
controlling what remains auditable. In the AI context, an operator
receiving a batch of deletion requests could strategically fulfil only
the requests corresponding to unfavorable decisions (evidence of bias),
while ignoring or delaying requests for benign entries. The resulting
log is legally compliant (GDPR-satisfied) but forensically distorted
(bias evidence erased). We formalise this attack in
Definition~\ref{def:sda}.

Existing solutions for redactable logs~\cite{ateniese2017,deuber2019}
address the mechanics of deletion but not the statistical distortion
it enables. They answer ``Can this entry be deleted?'' but not ``Does
this deletion distort the truth?'' This gap is not a limitation of any
specific scheme: the problem is absent from their threat models because
all prior work assumes the operator is a trusted party and the data
subject is the sole source of erasure requests.

This paper presents \textbf{Accountable Redaction}, a protocol that
enforces statistical integrity as a first-class condition on every log
transition. We model the log as a state machine: a redaction batch $R$
induces a transition $L_t \xrightarrow{R} L_{t+1}$, accepted if and
only if
\[
  \text{Valid}\!\left(L_t \xrightarrow{R} L_{t+1}\right) \;\iff\;
  \mathsf{Verify}(\pi_{\mathit{auth}},\, R)
  \;\land\;
  \mathsf{Verify}(\pi_{\mathit{stat}},\, L_t,\, L_{t+1},\, \varepsilon,\, \Pi)
\]
where $\pi_{\mathit{auth}}$ certifies that each erasure was requested by
the data subject or a legitimate authority, and $\pi_{\mathit{stat}}$ is
a ZK-SNARK proving that the conditional approval rates of all protected
groups shift by at most $\varepsilon$ after the batch. This ensures that
the right to privacy does not come at the expense of the public interest
in algorithmic accountability.

This work completes the accountability stack initiated by BTV:
\begin{itemize}
  \item \textbf{Paper~1 (Construction)~\cite{btv2026}:} A decision
  without evidence is a type error --- caught at compile time.
  \item \textbf{Paper~2 (Persistence)~\cite{btv-t-repo}:} A decision
  without persistence is a type error --- caught at delivery time.
  \item \textbf{Paper~3 (This paper):} A redaction without statistical
  justification is a verification failure --- caught at audit time.
\end{itemize}

We make three contributions:
\begin{enumerate}
  \item We define the Selective Deletion Attack
  (Definition~\ref{def:sda}) and formalise $\varepsilon$-statistical
  consistency (Definition~\ref{def:eps-stat}) as the security property
  that defeats it.
  \item We present the Accountable Redaction protocol
  (Section~\ref{sec:protocol}), combining Pedersen commitments over
  demographic attributes with a ZK-SNARK circuit for statistical
  verification, and prove it satisfies four security properties
  (Section~\ref{sec:security}).
  \item We provide a detailed gate-complexity analysis of the Noir
  circuit (Section~\ref{sec:zk}) and cost projections consistent with
  published Barretenberg benchmarks, demonstrating feasibility for
  production AI deployments (Section~\ref{sec:implementation}).
\end{enumerate}
